% Terjemahan Bahasa Indonesia dari chapter1.tex baris 926--1140.
% Sumber: Methods of Algebra, Volume 2, commit
% 9a5803ff77dd3257484cb177f851a73770a59dd3 (CC BY 4.0).
% stable-unit-id: o014.aljabr2.chapter1.kan-extensions

% segment-id: o014.li2.chapter1.kan-extensions.h001
\section{Ekstensi Kan}\label{sec:Kan-extension}
\hypertarget{o014.aljabr2.chapter1.kan-extensions}{ }

% segment-id: o014.li2.chapter1.kan-extensions.p001
Bagian ini bermula dari persoalan memperluas funktor. Tinjau kategori
$\mathcal{C}$, $\mathcal{D}$, $\mathcal{E}$ serta funktor $K$ dan $F$ seperti
dalam diagram berikut:
% segment-id: o014.li2.chapter1.kan-extensions.g001
\[\begin{tikzcd}
	\mathcal{C} \arrow[d, "K"'] \arrow[rd, bend left, "F"] & \\
	\mathcal{D} \arrow[dashed, r, "{\exists ? \; L}"'] & \mathcal{E}
\end{tikzcd}\]
Sesuai dengan semangat dasar teori kategori, kita ingin mencari funktor $L$
yang ditunjukkan oleh garis putus-putus, beserta isomorfisme alami
$F \simeq LK$. Persoalan ini pada umumnya tidak memiliki solusi; misalnya,
mungkin terdapat $f \in \Mor(\mathcal{C})$ sedemikian sehingga $Kf$ merupakan
isomorfisme, sedangkan $Ff$ bukan isomorfisme. Sebagai gantinya, setidaknya
kita dapat bertanya: adakah suatu hampiran terbaik? Jawabannya dicirikan oleh
sifat universal, dalam versi kiri dan kanan.

% segment-id: o014.li2.chapter1.kan-extensions.d001
\begin{definisi}[D.\ Kan]\label{def:Kan-extension}
	\index{ekstensi Kan (Kan extension)}
	\index[sym1]{LanKF@$\Lan_K F$}
	\index[sym1]{RanKF@$\Ran_K F$}
	Tinjau kategori $\mathcal{C}$, $\mathcal{D}$, $\mathcal{E}$, funktor
	$K: \mathcal{C} \to \mathcal{D}$ dan $F: \mathcal{C} \to \mathcal{E}$.
% segment-id: o014.li2.chapter1.kan-extensions.l001
	\begin{itemize}
% segment-id: o014.li2.chapter1.kan-extensions.i001
		\item \emph{Ekstensi Kan kiri} dari funktor $F$ di sepanjang $K$ adalah
		data $(\Lan_K F, \eta)$ berikut:
% segment-id: o014.li2.chapter1.kan-extensions.l002
		\begin{compactitem}
% segment-id: o014.li2.chapter1.kan-extensions.i002
			\item $\Lan_K F: \mathcal{D} \to \mathcal{E}$ merupakan funktor;
% segment-id: o014.li2.chapter1.kan-extensions.i003
			\item $\eta: F \to (\Lan_K F) K$ merupakan transformasi alami.
		\end{compactitem}
		Data ini harus memenuhi sifat universal berikut: untuk setiap data
		$L: \mathcal{D} \to \mathcal{E}$ dan $\xi: F \to LK$, terdapat tepat
		satu transformasi alami $\chi: \Lan_K F \to L$ sedemikian sehingga
		$\xi = (\chi K) \eta$ (komposisi vertikal/horizontal transformasi alami),
		atau, dalam diagram $2$-sel:
% segment-id: o014.li2.chapter1.kan-extensions.g002
		\[\begin{tikzcd}[column sep=large]
			\mathcal{C} \arrow[d, "K"'] \arrow[rd, bend left, "F", ""' {name=F}] & \\
			\mathcal{D} \arrow[r, "L"'] \arrow[Rightarrow, from=F, "\xi" description] & \mathcal{E}
% segment-id: o014.li2.chapter1.kan-extensions.g003
		\end{tikzcd} = \begin{tikzcd}[column sep=large]
			\mathcal{C} \arrow[d, "K"'] \arrow[rd, bend left, "F", ""' {name=F}] & \\
			\mathcal{D} \arrow[r, "{\Lan_K F}"' name=A] \arrow[Rightarrow, from=F, "\eta" description] \arrow[r, "L"{below, name=B}, rounded corners, to path={-- ([yshift=-2em]\tikztostart.south) -- ([yshift=-2em] \tikztotarget.south) [midway]\tikztonodes -- (\tikztotarget) }] & \mathcal{E} \arrow[Rightarrow, from=A, to=B, "\chi"]
		\end{tikzcd}. \]

% segment-id: o014.li2.chapter1.kan-extensions.i004
		\item \emph{Ekstensi Kan kanan} dari funktor $F$ di sepanjang $K$ adalah
		data $(\Ran_K F, \varepsilon)$ berikut:
% segment-id: o014.li2.chapter1.kan-extensions.l003
		\begin{compactitem}
% segment-id: o014.li2.chapter1.kan-extensions.i005
			\item $\Ran_K F: \mathcal{D} \to \mathcal{E}$ merupakan funktor;
% segment-id: o014.li2.chapter1.kan-extensions.i006
			\item $\varepsilon: (\Ran_K  F)K \to F$ merupakan transformasi alami.
		\end{compactitem}
		Data ini harus memenuhi sifat universal berikut: untuk setiap data
		$R: \mathcal{D} \to \mathcal{E}$ dan $\delta: RK \to F$, terdapat tepat
		satu transformasi alami $\theta: R \to \Ran_K F$ sedemikian sehingga
		$\delta = \varepsilon (\theta K)$, atau, dalam diagram $2$-sel:
% segment-id: o014.li2.chapter1.kan-extensions.g004
		\[\begin{tikzcd}[column sep=large]
			\mathcal{C} \arrow[d, "K"'] \arrow[rd, bend left, "F", ""' {name=F}] & \\
			\mathcal{D} \arrow[r, "R"'] \arrow[Rightarrow, to=F, "\delta" description] & \mathcal{E}
% segment-id: o014.li2.chapter1.kan-extensions.g005
		\end{tikzcd} = \begin{tikzcd}[column sep=large]
			\mathcal{C} \arrow[d, "K"'] \arrow[rd, bend left, "F", ""' {name=F}] & \\
			\mathcal{D} \arrow[r, "{\Ran_K F}"' name=A] \arrow[Rightarrow, to=F, "\varepsilon" description] \arrow[r, "R"{below, name=B}, rounded corners, to path={-- ([yshift=-2em]\tikztostart.south) -- ([yshift=-2em] \tikztotarget.south) [midway]\tikztonodes -- (\tikztotarget) }] & \mathcal{E} \arrow[Rightarrow, from=B, to=A, "\theta"]
		\end{tikzcd}. \]
	\end{itemize}
\end{definisi}

% segment-id: o014.li2.chapter1.kan-extensions.p002
Jika dalam definisi tersebut $\mathcal{C}$, $\mathcal{D}$, dan $\mathcal{E}$
diganti masing-masing dengan $\mathcal{C}^{\opp}$, $\mathcal{D}^{\opp}$, dan
$\mathcal{E}^{\opp}$, pola arah funktornya tetap sama, tetapi arah transformasi
alaminya berbalik. Oleh karena itu, ekstensi Kan kiri dan ekstensi Kan kanan
merupakan konsep yang saling dual.

% segment-id: o014.li2.chapter1.kan-extensions.p003
Dengan menggunakan kategori funktor $\mathcal{E}^{\mathcal{C}}$ dan
$\mathcal{E}^{\mathcal{D}}$, sifat universal ekstensi Kan kiri dan kanan
masing-masing menyatakan bijeksi
% segment-id: o014.li2.chapter1.kan-extensions.q001
\begin{equation}\label{eqn:Kan-univ}\begin{aligned}
	\Hom_{\mathcal{E}^{\mathcal{D}}}\left( \Lan_K F, L \right) & \xrightarrow{1:1} \Hom_{\mathcal{E}^{\mathcal{C}}}\left( F, LK \right) \\
	\chi & \longmapsto (\chi K)\eta , \\
	\Hom_{\mathcal{E}^{\mathcal{D}}}\left( R, \Ran_K F \right) & \xrightarrow{1:1} \Hom_{\mathcal{E}^{\mathcal{C}}}\left( RK, F \right) \\
	\theta & \longmapsto \varepsilon (\theta K);
\end{aligned}\end{equation}
inversnya masing-masing adalah pemetaan $(L, \xi) \mapsto \chi$ dan
$(R, \delta) \mapsto \theta$ dari Definisi \ref{def:Kan-extension}.

% segment-id: o014.li2.chapter1.kan-extensions.pr001
\begin{proposisi}\label{prop:Kan-uniqueness}
	Diberikan kategori $\mathcal{C}$, $\mathcal{D}$, $\mathcal{E}$ serta funktor
	$K: \mathcal{C} \to \mathcal{D}$ dan $F: \mathcal{C} \to \mathcal{E}$.
	Perkenalkan funktor tarik balik (yakni, funktor prapengomposisian)
	$K^*: \mathcal{E}^{\mathcal{D}} \to \mathcal{E}^{\mathcal{C}}$ di antara
	kategori-kategori funktor.
% segment-id: o014.li2.chapter1.kan-extensions.l004
	\begin{enumerate}[(i)]
% segment-id: o014.li2.chapter1.kan-extensions.i007
		\item Jika ada, ekstensi Kan kiri $(\Lan_K F, \eta)$ unik hingga
		isomorfisme unik di $\mathcal{E}^{\mathcal{D}}$; hal yang sama berlaku
		untuk ekstensi Kan kanan.
% segment-id: o014.li2.chapter1.kan-extensions.i008
		\item Jika ekstensi Kan kiri (atau ekstensi Kan kanan) di sepanjang $K$
		ada untuk setiap $F$, ekstensi-ekstensi tersebut dapat disusun menjadi
		funktor adjoin kiri bagi $K^*$,
		$\Lan_K: \mathcal{E}^{\mathcal{C}} \to \mathcal{E}^{\mathcal{D}}$
		(atau funktor adjoin kanan
		$\Ran_K: \mathcal{E}^{\mathcal{C}} \to \mathcal{E}^{\mathcal{D}}$);
		keluarga transformasi alami yang tersusun dari $\eta$ (atau
		$\varepsilon$) dalam Definisi \ref{def:Kan-extension} tepat merupakan
		unit (atau kounit) adjungsi tersebut.
	\end{enumerate}
\end{proposisi}
% segment-id: o014.li2.chapter1.kan-extensions.pf001
\begin{bukti}
	Untuk (i), cukup terapkan sifat universal dalam Definisi
	\ref{def:Kan-extension} dengan cara yang lazim.

	Untuk (ii), perhatikan bahwa $LK = K^*(L)$ dan $RK = K^*(R)$; relasi
	adjungsi yang dicari tepat merupakan isi \eqref{eqn:Kan-univ}. Pernyataan
	tentang unit atau kounit diperoleh dengan cara serupa; lihat
	\cite[Proposisi 2.6.5]{Li1}.
\end{bukti}

% segment-id: o014.li2.chapter1.kan-extensions.p004
Menurut konvensi, data $\eta$ atau $\varepsilon$ dalam ekstensi Kan sering kita
hilangkan dari notasi. Jika ekstensi Kan kiri (atau ekstensi Kan kanan) di
sepanjang $K$ ada untuk setiap $F$, funktor $\Lan_K$ (atau $\Ran_K$) yang
diperoleh dari (ii) juga memiliki sifat keunikan yang bersesuaian, berdasarkan
keunikan adjungsi \cite[Proposisi 2.6.10]{Li1}.

% segment-id: o014.li2.chapter1.kan-extensions.r001
\begin{catatan}
	Karena definisi ekstensi Kan hanya melibatkan komposisi $2$-sel, definisi
	tersebut dapat diperluas ke $2$-kategori umum; situasi di sini merupakan
	kasus khusus dalam $2$-kategori $\cate{Cat}$. Lihat
	\cite[\S 3.5]{Li1} untuk perinciannya.
\end{catatan}

% segment-id: o014.li2.chapter1.kan-extensions.p005
Untuk setiap kategori $\mathcal{C}$, tuliskan funktor tunggal
$\mathcal{C} \to \mathbf{1}$ sebagai $!$; menentukan funktor
$\mathbf{1} \to \mathcal{C}$ setara dengan menentukan sebuah objek di
$\mathcal{C}$. Untuk setiap kategori $\mathcal{E}$ dan objeknya $X$, tuliskan
$\Delta(X): \mathcal{C} \to \mathcal{E}$ untuk komposisi
$\mathcal{C} \xrightarrow{!} \mathbf{1} \xrightarrow{\text{konstan}\; X}
\mathcal{E}$; inilah funktor konstan yang memetakan setiap objek ke $X$ dan
setiap morfisme ke $\identity_X$.

% segment-id: o014.li2.chapter1.kan-extensions.e001
\begin{contoh}[Limit sebagai ekstensi Kan]\label{eg:limit-as-Kan}
	Misalkan $F: \mathcal{C} \to \mathcal{E}$ suatu funktor. Tinjau diagram
% segment-id: o014.li2.chapter1.kan-extensions.g006
	\[\begin{tikzcd}
		\mathcal{C} \arrow[d, "!"'] \arrow[rd, bend left, "F"] & \\
		\mathbf{1} & \mathcal{E}
	\end{tikzcd}\]
	Maka $\Lan_! F$ (atau $\Ran_! F$) ada jika dan hanya jika
	$\varinjlim F$ (atau $\varprojlim F$) ada di $\mathcal{E}$, dan dalam hal
	ini ekstensi Kan tersebut tidak lain adalah limit yang bersangkutan.

	Sebagai contoh, tinjau $(\Lan_! F, \eta)$. Funktor
	$\Lan_! F: \mathbf{1} \to \mathcal{E}$ dapat dipandang sebagai objek di
	$\mathcal{E}$, sedangkan $\eta$ merupakan transformasi alami
	$F \to \Delta(\Lan_! F)$ di $\mathcal{E}^{\mathcal{C}}$. Sifat universalnya
	setara dengan pernyataan bahwa data $(\Lan_! F, \eta)$ memberikan objek awal
	dalam kategori $(F/\Delta)$; inilah tepatnya uraian tentang $\varinjlim F$
	dalam \S\ref{sec:limit-functor}.
\end{contoh}

% segment-id: o014.li2.chapter1.kan-extensions.e002
\begin{contoh}[Adjungsi sebagai ekstensi Kan]
	Tinjau sepasang funktor:
% segment-id: o014.li2.chapter1.kan-extensions.g007
	\[\begin{tikzcd}
		F: \mathcal{C} \arrow[r, shift left] & \mathcal{D}: G \arrow[l, shift left]
	\end{tikzcd}\]
	Jika $(F, G)$ dilengkapi menjadi adjungsi dengan unit
	$\eta: \identity_{\mathcal{C}} \to GF$ dan kounit
	$\varepsilon: FG \to \identity_{\mathcal{D}}$, maka $(G, \eta)$ memberikan
	$\Lan_F (\identity_{\mathcal{C}})$ dan $(F, \varepsilon)$ memberikan
	$\Ran_G (\identity_{\mathcal{D}})$.

	Untuk menjelaskan hal ini, mula-mula kita tunjukkan bahwa funktor-funktor
	tarik balik berikut
% segment-id: o014.li2.chapter1.kan-extensions.g008
	\[\begin{tikzcd}
		G^*: \mathcal{C}^{\mathcal{C}} \arrow[r, shift left] &
		\mathcal{C}^{\mathcal{D}}: F^* \arrow[l, shift left]
	\end{tikzcd}\]
	juga dapat dilengkapi secara alami menjadi adjungsi. Perhatikan bahwa
	$\eta$ menginduksi
	$\identity_{\mathcal{C}}^* = \identity_{\mathcal{C}^{\mathcal{C}}}
	\to F^* G^* = (GF)^*$, dan secara serupa $\varepsilon$ menginduksi
	$G^* F^* \to \identity_{\mathcal{C}^{\mathcal{D}}}$; keduanya tetap kita
	tuliskan sebagai $\eta$ dan $\varepsilon$. Kita harus memeriksa identitas
	segitiga yang mencirikan adjungsi untuk
	$(G^*, F^*, \eta, \varepsilon)$. Secara ringkas, identitas-identitas ini
	merupakan konsekuensi formal dari penerapan tarik balik $(-)^*$ dalam
	kategori funktor pada identitas segitiga untuk
	$(F, G, \eta, \varepsilon)$; perinciannya diserahkan kepada pembaca.

	Adjungsi ini memberikan bijeksi kanonik
	$\Hom_{\mathcal{C}^{\mathcal{D}}}(\underbracket{G^* \identity_{\mathcal{C}}}_{= G}, L) \rightiso \Hom_{\mathcal{C}^{\mathcal{C}}}(\identity_{\mathcal{C}} , \underbracket{F^* L}_{=LF})$,
	dengan $L: \mathcal{D} \to \mathcal{C}$ sembarang; ketika $L = G$,
	$\identity_G$ dipetakan ke $\eta$. Jika sifat universal
	$\Lan_F (\identity_{\mathcal{C}})$ dituliskan sebagai bijeksi dalam
	\eqref{eqn:Kan-univ}, segera tampak bahwa $(G, \eta)$ memberikan
	$\Lan_F (\identity_{\mathcal{C}})$. Argumen untuk
	$\Ran_G (\identity_{\mathcal{D}})$ sepenuhnya serupa.\footnote{Catatan
	penerjemah: sumber menulis $\Ran_F(\identity_{\mathcal{C}})$ pada kalimat
	terakhir, tidak sesuai dengan pernyataan sebelumnya tentang
	$(F,\varepsilon)$. Target memperbaikinya menjadi
	$\Ran_G(\identity_{\mathcal{D}})$ (O014-C014).}
\end{contoh}

% segment-id: o014.li2.chapter1.kan-extensions.p006
Dalam praktiknya, konsep ekstensi Kan perlu diperkuat lebih lanjut dalam
sejumlah situasi\footnote{Salah satu penguatan yang lazim disebut ekstensi Kan
titik demi titik (pointwise Kan extension), yang tidak dibahas di sini.},
terutama dalam pembahasan funktor turunan nanti.

% segment-id: o014.li2.chapter1.kan-extensions.d002
\begin{definisi}\label{def:absolute-Kan}
	Tinjau kategori $\mathcal{C}$, $\mathcal{D}$, $\mathcal{E}$,
	$\mathcal{F}$ serta funktor $K: \mathcal{C} \to \mathcal{D}$ dan
	$F: \mathcal{C} \to \mathcal{E}$.\footnote{Catatan penerjemah: sumber tidak
	mencantumkan kategori $\mathcal{F}$ pada daftar awal, padahal kategori
	tersebut digunakan sebagai kodomain $M$. Target menambahkannya
	(O014-C015).}
% segment-id: o014.li2.chapter1.kan-extensions.l005
	\begin{itemize}
% segment-id: o014.li2.chapter1.kan-extensions.i009
		\item Misalkan ekstensi Kan kanan $(\Ran_K F, \varepsilon)$ ada. Kita
		katakan bahwa funktor $M: \mathcal{E} \to \mathcal{F}$
		\emph{mempertahankan} $\Ran_K F$ jika komposisi $2$-sel
% segment-id: o014.li2.chapter1.kan-extensions.g009
		\[\begin{tikzcd}[column sep=large]
			\mathcal{C} \arrow[d, "K"'] \arrow[rd, bend left, "F", ""' {name=F}] & & \\
			\mathcal{D} \arrow[r, "{\Ran_K F}"'] \arrow[Rightarrow, to=F, "\varepsilon" description] & \mathcal{E} \arrow[r, "M"'] & \mathcal{F}
% segment-id: o014.li2.chapter1.kan-extensions.g010
		\end{tikzcd} =: \begin{tikzcd}
			\mathcal{C} \arrow[rr, "MF", ""{name=MF}] \arrow[rd, "K"'] & & \mathcal{F}  \\
			& \mathcal{D} \arrow[ru, "{M \Ran_K F}"'] \arrow[Rightarrow, to=MF, "M\varepsilon" description] &
		\end{tikzcd}\]
		memberikan ekstensi Kan kanan dari $MF$ di sepanjang $K$.
% segment-id: o014.li2.chapter1.kan-extensions.i010
		\item Misalkan ekstensi Kan kiri $(\Lan_K F, \eta)$ ada. Kita katakan
		bahwa funktor $M: \mathcal{E} \to \mathcal{F}$
		\emph{mempertahankan} $\Lan_K F$ jika komposisi $2$-sel
% segment-id: o014.li2.chapter1.kan-extensions.g011
		\[\begin{tikzcd}[column sep=large]
			\mathcal{C} \arrow[d, "K"'] \arrow[rd, bend left, "F", ""' {name=F}] & & \\
			\mathcal{D} \arrow[r, "{\Lan_K F}"'] \arrow[Rightarrow, from=F, "\eta" description] & \mathcal{E} \arrow[r, "M"'] & \mathcal{F}
% segment-id: o014.li2.chapter1.kan-extensions.g012
		\end{tikzcd} =: \begin{tikzcd}
			\mathcal{C} \arrow[rr, "MF", ""{name=MF}] \arrow[rd, "K"'] & & \mathcal{F}  \\
			& \mathcal{D} \arrow[ru, "{M \Lan_K F}"'] \arrow[Rightarrow, from=MF, "M\eta" description] &
		\end{tikzcd}\]
		memberikan ekstensi Kan kiri dari $MF$ di sepanjang $K$.
% segment-id: o014.li2.chapter1.kan-extensions.i011
		\item Jika ekstensi Kan kiri $(\Lan_K F, \eta)$ (atau ekstensi Kan kanan
		$(\Ran_K F, \varepsilon)$) dipertahankan oleh setiap funktor yang
		berdomain $\mathcal{E}$, ekstensi Kan tersebut disebut
		\emph{absolut}.
		\index{ekstensi Kan!absolut (absolute)}
	\end{itemize}
\end{definisi}

% segment-id: o014.li2.chapter1.kan-extensions.p007
Langkah berikutnya adalah menetapkan hubungan antara adjungsi dan ekstensi Kan
absolut. Hasil ini sangat berguna dalam kajian lanjut tentang funktor turunan;
mempelajari buktinya juga membantu membiasakan diri dengan ekstensi Kan dan
manipulasi $2$-sel. Tinjau sepasang funktor adjoin
% segment-id: o014.li2.chapter1.kan-extensions.g013
$\begin{tikzcd}
	F: \mathcal{C} \arrow[r, shift left] & \mathcal{C}' \arrow[l, shift left] : G
\end{tikzcd}$,
yang ditentukan oleh unit $\eta: \identity_{\mathcal{C}} \to GF$ dan kounit
$\varepsilon: FG \to \identity_{\mathcal{C}'}$. Diberikan pula funktor
% segment-id: o014.li2.chapter1.kan-extensions.q002
\[ K: \mathcal{C} \to \mathcal{D}, \quad K': \mathcal{C}' \to \mathcal{D}'. \]
Selanjutnya kita akan menggunakan diagram $2$-sel secara sistematis untuk
menyatakan funktor, transformasi alami di antaranya, dan komposisinya, agar
argumen menjadi lebih ringkas. Dalam manipulasi $2$-sel, kita akan menukar
komposisi vertikal dan horizontal tanpa penjelasan lebih lanjut; tindakan ini
sah, sebagaimana dijelaskan dalam \cite[Lema 2.2.7]{Li1}.

% segment-id: o014.li2.chapter1.kan-extensions.th001
\begin{teorema}[{G.\ Maltsiniotis \cite{Ma07}}]\label{prop:Quillen-adjunction-gen}
	Misalkan $K'F$ mempunyai ekstensi Kan kanan absolut di sepanjang $K$,
	yang dinotasikan dengan $\mathrm{L}F$, sedangkan $KG$ mempunyai ekstensi
	Kan kiri absolut di sepanjang $K'$, yang dinotasikan dengan $\mathrm{R}G$.
	Data terkait ditampilkan dalam diagram berikut:
% segment-id: o014.li2.chapter1.kan-extensions.g014
	\[\begin{tikzcd}
		\mathcal{C} \arrow[r, "F"] \arrow[d, "K"'] & \mathcal{C}' \arrow[d, "{K'}"] \\
		\mathcal{D} \arrow[r, "{\mathrm{L}F}"'] \arrow[Rightarrow, ru, "\alpha" description] & \mathcal{D}'
% segment-id: o014.li2.chapter1.kan-extensions.g015
	\end{tikzcd} \quad \begin{tikzcd}
		\mathcal{C}' \arrow[r, "G"] \arrow[d, "{K'}"'] & \mathcal{C} \arrow[d, "K"] \arrow[Rightarrow, ld, "\beta" description] \\
		\mathcal{D}' \arrow[r, "{\mathrm{R}G}"'] & \mathcal{D} .
	\end{tikzcd} \]
	Dalam keadaan ini, terdapat tepat satu
	$\underline{\eta}: \identity_{\mathcal{D}} \to \mathrm{R}G \circ \mathrm{L}F$
	dan tepat satu
	$\underline{\varepsilon}: \mathrm{L}F \circ \mathrm{R}G \to
	\identity_{\mathcal{D}'}$ sedemikian sehingga persamaan komposisi $2$-sel
	berikut berlaku:
% segment-id: o014.li2.chapter1.kan-extensions.q003
	\begin{equation*}\begin{gathered}
% segment-id: o014.li2.chapter1.kan-extensions.g016
		\begin{tikzcd}[every cell/.append style = {font = \small}]
			\mathcal{C} \arrow[r] \arrow[rr, bend left=75, "{\identity_{\mathcal{C}}}", ""' {name=I}] & \mathcal{C}' \arrow[r] \arrow[d] \arrow[Rightarrow, from=I, "\eta" description] & \mathcal{C} \arrow[d] \arrow[Rightarrow, ld, "\beta" description] \\
			& \mathcal{D}' \arrow[r, "{\mathrm{R}G}"'] & \mathcal{D}
% segment-id: o014.li2.chapter1.kan-extensions.g017
		\end{tikzcd} = \begin{tikzcd}[every cell/.append style = {font = \small}]
			\mathcal{C} \arrow[r] \arrow[d] & \mathcal{C}' \arrow[d] & \\
			\mathcal{D} \arrow[r, "{\mathrm{L}F}"'] \arrow[rr, bend right=75, "{\identity_{\mathcal{D}}}"', "" {name=I}] \arrow[Rightarrow, ru, "\alpha" description] & \mathcal{D}' \arrow[r, "{\mathrm{R}G}"'] \arrow[Rightarrow, from=I, "\underline{\eta}" description] & \mathcal{D}
		\end{tikzcd} \\
% segment-id: o014.li2.chapter1.kan-extensions.g018
		\begin{tikzcd}[every cell/.append style = {font = \small}]
			\mathcal{C}' \arrow[r] \arrow[bend left=75, rr, "{\identity_{\mathcal{C}'}}", ""' {name=I}] & \mathcal{C} \arrow[r] \arrow[d] \arrow[Rightarrow, to=I, "\varepsilon" description] & \mathcal{C}' \arrow[d] \\
			& \mathcal{D} \arrow[r, "{\mathrm{L}F}"'] \arrow[Rightarrow, ru, "\alpha" description] & \mathcal{D}'
% segment-id: o014.li2.chapter1.kan-extensions.g019
		\end{tikzcd} = \begin{tikzcd}
			\mathcal{C}' \arrow[r] \arrow[d] & \mathcal{C} \arrow[d] \arrow[Rightarrow, ld, "\beta" description] & \\
			\mathcal{D}' \arrow[r, "{\mathrm{R}G}"'] \arrow[rr, bend right=75, "{\identity_{\mathcal{D}'}}"', "" {name=I}] & \mathcal{D} \arrow[r, "{\mathrm{L}F}"'] \arrow[Rightarrow, to=I, "\underline{\varepsilon}" description] & \mathcal{D}'
		\end{tikzcd}
	\end{gathered}\end{equation*}
	Selanjutnya, $\underline{\eta}$ dan $\underline{\varepsilon}$ tersebut
	membuat $(\mathrm{L}F, \mathrm{R}G)$ menjadi suatu adjungsi.
\end{teorema}
% segment-id: o014.li2.chapter1.kan-extensions.pf002
\begin{bukti}
	%%%%%%%%%%%%%%%%%%%%%%%%%
	Tugas pertama adalah membuktikan keberadaan dan keunikan
	$\underline{\eta}$ serta $\underline{\varepsilon}$. Karena $\mathrm{R}G$
	merupakan ekstensi Kan kiri absolut, $2$-sel
% segment-id: o014.li2.chapter1.kan-extensions.g020
	\[\begin{tikzcd}
		\mathcal{C}' \arrow[r, "G"] \arrow[d, "{K'}"'] & \mathcal{C} \arrow[d, "K"] \arrow[Rightarrow, ld, "\beta" description] & \\
		\mathcal{D}' \arrow[r, "{\mathrm{R}G}"'] & \mathcal{D} \arrow[r, "{\mathrm{L}F}"'] & \mathcal{D}'
	\end{tikzcd} = (\mathrm{L}F)\beta : \mathrm{L}F \circ K \circ G \to \mathrm{L}F \circ \mathrm{R}G \circ K' \]
	membuat $\mathrm{L}F \circ \mathrm{R}G$ menjadi ekstensi Kan kiri dari
	$\mathrm{L}F \circ K \circ G$ di sepanjang $K'$. Di sisi lain, kita juga
	mempunyai
% segment-id: o014.li2.chapter1.kan-extensions.g021
	\[\begin{tikzcd}[every cell/.append style = {font = \small}]
		\mathcal{C}' \arrow[r, "G"] \arrow[bend left=75, rr, "{\identity_{\mathcal{C}'}}", ""' {name=I}] & \mathcal{C} \arrow[r, "F"] \arrow[d, "K"'] \arrow[Rightarrow, to=I, "\varepsilon" description] & \mathcal{C}' \arrow[d, "{K'}"] \\
		& \mathcal{D} \arrow[r, "{\mathrm{L}F}"'] \arrow[Rightarrow, ru, "\alpha" description] & \mathcal{D}'
	\end{tikzcd} = \varepsilon(\alpha G) : \mathrm{L}F \circ K \circ G \to K' \circ \identity_{\mathcal{C}'} = K' . \]
	Dengan demikian, sifat universal ekstensi Kan kiri menentukan tepat satu
	$\underline{\varepsilon}$ yang memenuhi persamaan komposisi $2$-sel dalam
	pernyataan teorema.
	%%%%%%%%%%%%%%%%%%%%%%%%%

	Kasus $\underline{\eta}$ bersifat dual. Cukup perhatikan bahwa
	$(\mathrm{R}G)\alpha:
	\mathrm{R}G \circ \mathrm{L}F \circ K \to
	\mathrm{R}G \circ K' \circ F$ juga membuat
	$\mathrm{R}G \circ \mathrm{L}F$ menjadi ekstensi Kan kanan, sebab
	$\mathrm{L}F$ merupakan ekstensi Kan kanan absolut.

	Tinggal memeriksa identitas segitiga bagi unit $\underline{\eta}$ dan kounit
	$\underline{\varepsilon}$, yakni persamaan komposisi $2$-sel berikut:
% segment-id: o014.li2.chapter1.kan-extensions.q004
	\begin{equation*}\begin{gathered}
% segment-id: o014.li2.chapter1.kan-extensions.g022
		\begin{tikzcd}[column sep=small, every cell/.append style = {font = \small}]
			\mathcal{D}' \arrow[r] \arrow[bend right=60, rr, "\identity"', "" {name=I}] & \mathcal{D} \arrow[r] \arrow[rr, bend left=60, "\identity", "" {name=J}] \arrow[Rightarrow, to=I] & \mathcal{D}' \arrow[r] \arrow[Rightarrow, from=J] & \mathcal{D}
% segment-id: o014.li2.chapter1.kan-extensions.g023
		\end{tikzcd} = \begin{tikzcd}[every cell/.append style = {font = \small}]
			\mathcal{D}' \arrow[r, bend left=60, "{\mathrm{R}G}", ""' {name=A}] \arrow[r, bend right=60, "{\mathrm{R}G}"', "" {name=B}] & \mathcal{D} \arrow[Rightarrow, from=A, to=B, "\identity" description]
		\end{tikzcd} , \quad
% segment-id: o014.li2.chapter1.kan-extensions.g024
		\begin{tikzcd}[column sep=small, every cell/.append style = {font = \small}]
			\mathcal{D} \arrow[r] \arrow[bend left=60, rr, "\identity", ""' {name=I}] & \mathcal{D}' \arrow[r] \arrow[Rightarrow, from=I] \arrow[bend right=60, rr, "\identity"', "" {name=J}] & \mathcal{D} \arrow[r] \arrow[Rightarrow, to=J] & \mathcal{D}'
% segment-id: o014.li2.chapter1.kan-extensions.g025
		\end{tikzcd} = \begin{tikzcd}[every cell/.append style = {font = \small}]
			\mathcal{D} \arrow[r, bend left=60, "{\mathrm{L}F}", ""' {name=A}] \arrow[r, bend right=60, "{\mathrm{L}F}"', "" {name=B}] & \mathcal{D}' \arrow[Rightarrow, from=A, to=B, "\identity" description]
		\end{tikzcd},
	\end{gathered}\end{equation*}
	dengan label pada anak panah yang tidak menimbulkan ambiguitas dihilangkan.
	Kita buktikan persamaan pertama. Berdasarkan sifat universal
	$(\mathrm{R}G, \beta)$ sebagai ekstensi Kan kiri, cukup ditunjukkan bahwa
	$2$-sel yang diperoleh dengan mengomposisikan kedua ruas dengan $\beta$
	adalah sama. Untuk ruas kiri persamaan pertama, mulailah dari
% segment-id: o014.li2.chapter1.kan-extensions.g026
	\[\begin{tikzcd}[every cell/.append style = {font = \small}]
		\mathcal{C}' \arrow[r] \arrow[d] & \mathcal{C} \arrow[d] \arrow[Rightarrow, ld] & & \\
		\mathcal{D}' \arrow[r] \arrow[bend right=60, rr, "\identity"', "" {name=I}] & \mathcal{D} \arrow[r] \arrow[rr, bend left=60, "\identity", "" {name=J}] \arrow[Rightarrow, to=I, "\underline{\varepsilon}"] & \mathcal{D}' \arrow[r] \arrow[Rightarrow, from=J, "\underline{\eta}"] & \mathcal{D}
% segment-id: o014.li2.chapter1.kan-extensions.g027
	\end{tikzcd} = \begin{tikzcd}[every cell/.append style = {font = \small}]
		\mathcal{C}' \arrow[r] \arrow[rr, bend left=40, "\identity", ""' {name=T}] & \mathcal{C} \arrow[r] \arrow[d] \arrow[Rightarrow, to=T, "\varepsilon" inner sep=0.4em] & \mathcal{C}' \arrow[d] & \\
 		& \mathcal{D} \arrow[r] \arrow[rr, bend right=60, "\identity"', "" {name=I}] \arrow[Rightarrow, ru] & \mathcal{D}' \arrow[r] \arrow[Rightarrow, from=I, "\underline{\eta}"'] & \mathcal{D}
	\end{tikzcd}\]
	di mana kita menggunakan pencirian $\underline{\varepsilon}$. Selanjutnya,
	dengan menggunakan pencirian $\underline{\eta}$ serta identitas segitiga
	yang dipenuhi $\eta$ dan $\varepsilon$, ruas kanan di atas menjadi
% segment-id: o014.li2.chapter1.kan-extensions.g028
	\[\begin{tikzcd}[every cell/.append style = {font = \small}]
		\mathcal{C}' \arrow[r] \arrow[rr, bend right=75, "\identity"', "" {name=A}] & \mathcal{C} \arrow[r] \arrow[rr, bend left=75, "\identity", ""' {name=B}] \arrow[Rightarrow, to=A, "\varepsilon"] & \mathcal{C}' \arrow[r] \arrow[d] \arrow[Rightarrow, from=B, "\eta"] & \mathcal{C} \arrow[d] \arrow[Rightarrow, ld, "\beta" description] \\
		& & \mathcal{D}' \arrow[r] & \mathcal{D}
% segment-id: o014.li2.chapter1.kan-extensions.g029
	\end{tikzcd} = \begin{tikzcd}[column sep=large, every cell/.append style = {font = \small}]
		\mathcal{C}' \arrow[r] \arrow[bend left=75, r, ""' {name=T}] \arrow[r, "" {name=B}] \arrow[d] & \mathcal{C} \arrow[d] \arrow[Rightarrow, ld, "\beta" description] \\
		\mathcal{D}' \arrow[r] & \mathcal{D} \arrow[Rightarrow, from=T, to=B, "\identity"]
% segment-id: o014.li2.chapter1.kan-extensions.g030
	\end{tikzcd} = \begin{tikzcd}[column sep=large, every cell/.append style = {font = \small}]
		\mathcal{C}' \arrow[r] \arrow[d] & \mathcal{C} \arrow[d]  \arrow[Rightarrow, ld, "\beta" description] \\
		\mathcal{D}' \arrow[r, ""' {name=T}] \arrow[r, bend right=75, "" {name=B}] & \mathcal{D} \arrow[Rightarrow, from=T, to=B, "\identity"]
	\end{tikzcd}\]
	Inilah persamaan segitiga pertama. Pembuktian persamaan kedua bersifat dual.
\end{bukti}
